\chapter*{Введение}
\addcontentsline{toc}{chapter}{Введение}

Современная электроэнергетика опирается на большое количество измерительных, расчетных и эксплуатационных данных.
Даже относительно простые инженерные задачи --- анализ режима, проверка качества данных, оценка состояния оборудования, краткосрочный прогноз нагрузки или поиск аномалий --- сегодня все чаще требуют систематической обработки табличной и временной информации.

В этих условиях возникает потребность не просто в изучении языка программирования или отдельных библиотек, а в освоении цельного рабочего контура, в котором данные, вычисления, визуализация, интерпретация и фиксация результатов объединены в единую воспроизводимую среду.

Настоящее пособие строится именно в такой логике.
В качестве основной рабочей среды используется формат Jupyter-блокнотов, а содержательная часть курса ориентирована на прикладной анализ данных в задачах электроэнергетики.
При этом акцент делается не на интерфейсных особенностях конкретной оболочки, а на последовательности инженерного исследования: описание задачи, получение и проверка данных, предобработка, формирование признаков, выбор метрик, построение базовых и более содержательных моделей, интерпретация результатов и подготовка итогового вывода.

Материал организован по принципу перехода от общей постановки к практическим кейсам.
Сначала рассматриваются среда работы, источники данных и первичный контроль качества, затем --- разведочный анализ, формирование признаков, выбор схемы оценки качества и базовые модели.
После этого вводятся три сквозных кейса: классификация, регрессия и временные ряды.

Каждая глава сопровождается связанным Jupyter-блокнотом.
Такое решение позволяет развести два уровня представления материала:
\begin{itemize}[leftmargin=1.25cm]
\item в тексте пособия фиксируются понятия, логика действий, формулы, ограничения и инженерная интерпретация;
\item в блокнотах воспроизводится полный вычислительный сценарий с кодом, таблицами, графиками и результатами.
\end{itemize}

Подобное разделение особенно важно для учебной практики.
С одной стороны, оно не перегружает основное изложение длинными листингами.
С другой стороны, оно позволяет сохранить исполнимость примеров и обеспечить повторяемость вычислений.

Пособие адресовано студентам, магистрантам, преподавателям и инженерам, осваивающим методы анализа данных на табличных и временных рядах в электроэнергетике.
